\chapter{A Génese Dependente III: As Formações do Eu}

% REVISÃO (ortografia): capítulo escrito segundo o Acordo Ortográfico de 1990 ("perceção", "afetado", "refletir", "caráter", "direcionada", "objetos", "ações", "ato"); os capítulos 00–19 usam a grafia pré-Acordo. Uniformizar em todo o livro.
% REVISÃO (terminologia): "formação kármica" / "formações kármicas" — noutros pontos do livro usa-se "kammica" (de "kamma"); uniformizar. E "Buda" (linhas 80, 96, 161, 163) → "Buddha".
\emph{Avijjāpaccayā saṅkhārā:} Isto significa: «a ignorância condiciona a
formação kármica», ou seja, o corpo e a mente tal como definidos pelos cinco
\emph{khandhas.} É quando agimos a partir de uma posição de ignorância, sem
compreender a verdade, e tudo o que experimentamos, fazemos, dizemos e sentimos
é condicionado por essa ignorância. Absolutamente tudo.

É aqui que a visão do «eu» constitui um verdadeiro ponto cego. Quando pensamos
na formação kármica como «eu» em vez de como «não-eu», então tudo o que
acontece, tudo o que é vivido, é atribuído a essa noção de pessoa -- o «eu»
enquanto perceção. Isto é \emph{avijjāpaccayā saṅkhārā.}

Se tivermos a percepção de que todas as condições são impermanentes, que todo o
dhamma é não-eu, então há conhecimento ou \emph{vijjā}, e verdade ou Dhamma, em
vez de ignorância (\emph{avijjā}) e kamma habitual \emph{(saṅkhāra).} Há o
conhecimento do Dhamma, a verdade de como as coisas são. Então, tudo o resto
segue o mesmo caminho, tudo é visto tal como é. Não há distorção: a consciência,
os cinco agregados e o mundo sensorial são vistos como Dhamma, em vez de como
eu.

Afinal, qual é o teu sofrimento na vida? Por que sofres? Se investigares, podes
sempre remontar a \emph{avijjāpaccayā saṅkhārā.} Há sofrimento natural: passar
fome, envelhecer e adoecer, mas tudo isso é suportável. Não é nada que não
possamos suportar. A doença, a velhice e a morte são coisas com as quais podemos
sempre lidar. Isso não é sofrimento real. Mas o sofrimento é a ganância, o ódio
e a ilusão que produzimos através da visão do eu, ao levarmos tudo a peito. As
criações e os apegos a visões erradas, preconceitos, parcialidades e todos os
horrores pelos quais somos responsáveis podem todos ser atribuídos a
\emph{avijjāpaccayā saṅkhārā}.

Não podemos realmente esperar grande melhoria se continuarmos a insistir em ser
ignorantes, presos à visão do eu. Mesmo que possamos melhorar ligeiramente as
condições ao tentar ser uma boa pessoa, enquanto houver apego à visão do eu, há
ilusão; por isso, até a bondade que praticamos provém da ilusão. Isso não conduz
ninguém para fora do sofrimento. Se não tivermos sabedoria, então muitas vezes
tentamos fazer o bem, mas acabamos por causar todo o tipo de problemas, enquanto
pensamos que podemos dizer aos outros o que é bom para eles.

Como é que as coisas são neste momento? O teu corpo está sentado, não está?
Podes sentir coisas -- prazer, dor, calor ou frio ou o que quer que seja. É
assim que as coisas são. Não há «eu» nisso; não estamos a criar o «eu». Quando
voltamos a nossa atenção para a forma como as coisas são, podemos ver o que
fazemos quando criamos o «eu» e o «meu» no momento: o que penso, o que sinto, o
que quero, o que não quero, o que gosto, o que não gosto, ou podemos estar
conscientes dos «eus» que criamos nos outros: as minhas opiniões sobre ti. Sofri
muito por criar pessoas na minha mente, não porque alguém fosse realmente cruel
comigo, mas por causa de todas as coisas que costumava inventar sobre mim e
sobre os outros -- o medo do que os outros pensavam; as invejas, a cobiça, a
possessividade. Eu tinha os meus preconceitos e opiniões sobre as pessoas, o que
achava que elas realmente estavam a tramar, e as minhas suspeitas sobre o que
elas realmente queriam. Assim, esse sofrimento surge da imagem que criamos de
nós próprios e dos outros, dos nossos pais e das pessoas mais próximas de nós.

O que é o sofrimento? Pergunta-te a sério: o que é o sofrimento da tua vida?
Ontem, os ventos frios sopravam através de mim enquanto caminhava ali no campo.
Isso é sofrimento? Eu podia transformá-lo em sofrimento. «Odeio este vento frio
e não gosto dele.» Mas, na verdade, estava tudo bem. Quero dizer, era algo que
eu conseguia suportar perfeitamente. Se não criasse nada na minha mente a esse
respeito, era apenas vento frio -- nada mais.

No entanto, podemos passar o tempo em Amaravati a criar atitudes em relação a
monges, monjas e leigos. É possível transformar monjas seniores em verdadeiros
ogros, não é? Podemos ter opiniões fortes sobre a senioridade. Se estivermos
numa posição de senioridade, podemos ficar muito apegados a ela. «Sou mais
antiga do que tu. Tu és apenas um monge novato. Faz isso. Eu sou o chefe.»
Assim, podemos criar a imagem de nós próprios como monges seniores. Mas não
estamos aqui para criar kamma com base na ignorância. As convenções que temos
são meramente meios expedientes. São simplificações, acordos morais e acordos
comunitários, para tornar a vida simples e descomplicada e também para nos
permitir refletir sobre a forma como nos relacionamos com as pessoas; pessoas
seniores, pessoas do mesmo nível, pessoas juniores.

% FIXME: Verse reference is not 860
% https://www.digitalpalireader.online/_dprhtml/index.html?loc=k.4.0.0.3.0.9.m
% This verse?
% Vītagedho amaccharī, Na ussesu vadate muni; Na samesu na omesu, Kappaṁ neti akappiyo.
% \quoteRef{\href{https://suttacentral.net/snp4.10/pli/ms}{Snp 4.10 Purābheda Sutta}}

O Buda disse\footnote{Por exemplo, Sutta-Nipata 860: «Tendo acabado com a
  inveja e a ganância, o sábio não fala de ``superior'', nem de ``inferior'',
  nem de ``igual''.»}: «A visão de que todos são iguais é uma ilusão.» «Eu sou»
é uma ilusão, se essa identidade se basear na ignorância. Mas quando há
\emph{vijjā}, então «eu sou» é meramente uma realidade convencional. É apenas a
forma como falamos: «Tenho fome», ou «Sou Sumedho Bhikkhu», mas não é uma
pessoa.

Quando há \emph{avijjā}, isso condiciona os \emph{saṅkhārās} que condiciona a
consciência ou \emph{viññāṇa}. A consciência condiciona a mentalidade e a
corporeidade (\emph{nāma-rūpa}), que condicionam as bases sensoriais
(\emph{saḷāyatana}), que condicionam o contacto (\emph{phassa}), que condiciona
a sensação (\emph{vedanā}). Quando a ignorância é a condição primária, o resto é
todo afetado por isso. O mundo sensorial, o corpo e a mente, estão relacionados
com ela em termos de «eu» e «meu». Esta é a visão do eu. Agora, em contraste com
o Brahmanismo, onde os hindus falam do Atman ou do «eu superior», o Um, quando o
Buda fala do eu, este está relacionado com o apego aos cinco \emph{khandhas}, ao
corpo, à sensação, às perceções, às volições e à consciência; o apego a isso, a
ignorância, condiciona as formações kármicas. Tudo isto cria uma sensação de um
eu.

Esta visão do eu começa a tornar-se forte por volta dos sete anos de idade.
Vais para a escola, competes e comparas, e esta forte sensação de um eu começa a
ser condicionada na tua mente. Lembro-me de os primeiros cinco ou seis anos da
minha vida terem sido mágicos e, depois dos seis ou sete, ter começado a piorar
cada vez mais. Antes disso, não havia grande sensação de eu.

Num país como os Estados Unidos -- que, na verdade, é um país muito agradável
--, há uma ênfase na visão do eu. Não há uma grande quantidade de sabedoria
nesse país e a visão da personalidade é, sem dúvida, o tema dominante. «Sou um
indivíduo. Tenho os meus direitos. Posso fazer o que quiser. Não me podes dizer
o que fazer. Quem te achas que és? Sou tão bom quanto tu. Deixa-me em paz.» Os
americanos igualitários têm um forte apego a serem indivíduos com uma
personalidade fascinante, um caráter autêntico, um «bom rapaz». Esta é a ênfase
americana no nível pessoal. Ser um «bom rapaz» é aceitável; não há nada de
errado nisso, mas, como \emph{avijjāpaccayā saṅkhārā}, isso só pode trazer
sofrimento. Quando há ignorância e visão do eu, o «bom rapaz» vai sempre sofrer.

\emph{Avijjā} condiciona os \emph{saṅkhārās}, que condiciona a consciência, que
condiciona a mentalidade-corporeidade, que condiciona os seis sentidos, que
condicionam o contacto e, em seguida, a sensação, e então a sensação condiciona
o desejo -- a ligação \emph{vedanā-taṇhā}. Podes perceber que, se estiveres
preso ao apego à visão da personalidade ou ao eu, então haverá desejo, apego
(\emph{upādāna}) e devir (\emph{bhava}). Ficarás perdido nesse padrão porque,
quando há ignorância \emph{neste} momento, isso afeta tudo -- a consciência, os
sentidos e os objetos dos sentidos, a sensação e, em seguida, o desejo entra em
% REVISÃO: aspas/frase corrompidas — "«Eu quero livrar-me do -- «Eu quero.»": a primeira aspa não fecha e "livrar-me do" fica incompleto. Corrigir, p. ex.: «Eu quero livrar-me disto.» «Eu quero.»
cena. «Eu quero algo.» «Eu quero ser feliz.» «Eu quero tornar-me.» «Eu quero
livrar-me do -- «Eu quero.»

Examine o desejo durante este retiro; procure realmente compreender o que é o
desejo. A partir da minha própria reflexão sobre o assunto, vejo que é sempre
energia direcionada para algo, quer seja inquieta e dispersa, quer esteja
direcionada para algo definido. Existe um forte desejo de nos livrarmos das
coisas de que não gostamos o mais rapidamente possível. Queremos obter o que
queremos instantaneamente e livrar-nos do que não gostamos instantaneamente; já
não valorizamos a paciência na nossa sociedade. Queremos eficiência. Tudo parece
bem e, de repente, surge algo que faz confusão e temos de limpar imediatamente,
porque não queremos obstáculos, impedimentos ou qualquer coisa desagradável.
Queremos livrar-nos disso rapidamente, por isso somos muito impacientes e
podemos ficar muito chateados e irritados com as coisas devido a este desejo de
nos livrarmos, este \emph{vibhava-taṇhā}.

O desejo de nos tornarmos, a ambição, o \emph{bhava-taṇhā}, é frequentemente
uma motivação na vida religiosa -- queremos tornar-nos uma personalidade
iluminada. Por isso, o \emph{bhava-taṇhā} e o \emph{vibhava-taṇhā} devem ser
estudados e examinados.

Podes refletir sobre eles; podes ouvir estes desejos: «Quero alcançar a
iluminação.» «Quero alcançar o \emph{samādhi}. Quero tirar o máximo partido deste
retiro para que possa obter algum tipo de realização ou conquista com tudo
isto.» Também queremos livrar-nos de coisas: «Espero livrar-me de toda a minha
luxúria e raiva durante este retiro. Espero livrar-me da inveja para nunca mais
ter de sentir inveja. Durante este retiro estou a trabalhar na inveja. Estou a
trabalhar na dúvida ou no medo -- se conseguir livrar-me do meu medo até ao
final deste retiro, não terei mais medo, porque vou entrar lá e aniquilar o
medo.» Isso é \emph{vibhava-taṇhā}! «Há algo de errado comigo e tenho de
corrigir isso. Tenho de me tornar outra coisa, livrando-me destas coisas más,
destas coisas erradas em mim.» São todos os «eu sou» e os «eu e o meu».

\emph{Kāma-taṇhā} é bastante óbvio -- é o desejo por experiências sensoriais
prazerosas. Estas formas de desejo devem ser conhecidas e \emph{compreendidas.}
A armadilha é que tendemos a pensar que o Buda nos ensina a livrar-nos dos
nossos desejos. É assim que algumas pessoas interpretam o budismo. Mas isso está
errado: o Buda ensinou-nos a \emph{observar} e \emph{compreender} o desejo para
que não nos apeguemos a ele! Isso não significa livrar-nos do desejo, mas
compreendê-lo verdadeiramente para que o desejo não nos possa mais iludir. O
desejo de nos livrarmos do desejo continua a ser desejo; não está a observar o
desejo. Com esse desejo, estás apenas a agarrar-te à perceção de que não deves
ter desejos e que tens de te livrar deles. Mas, ao compreender a Originação
Dependente, vemos o \emph{taṇhā} como Dhamma e não como o eu -- estás a olhar
para o \emph{taṇhā}, o desejo, como aquilo que surge e cessa. Isso é Dhamma, não é? Em
vinte e cinco anos de olhar atento e observação minuciosa, não encontrei um
único desejo que surja e continue a surgir. Se algum de vocês encontrar um, por
favor, diga-me.

\emph{Kāma-taṇhā} é bastante grosseiro e bastante óbvio -- quero algo para
comer, ou desejos sexuais. Mas \emph{vibhava-taṇhā} pode ser muito subtil,
justo e importante. E pode-se ser iludido por essa qualidade de justiça. O
desejo de se livrar do mal pode parecer tão certo. Podemos dedicar as nossas
vidas a livrar-nos dos males deste mundo e tornar-nos fanáticos. É isto que se
pode ver nos problemas sociais modernos. Existem as tendências degeneradas desta
sociedade que se traduzem em aberrações sexuais e drogas e, depois, existem as
formas muito justas dos fundamentalistas que condenam o comportamento
degenerado, dissoluto e imoral de um elemento da sociedade. Mas estamos a olhar
para o próprio desejo, desde as formas grosseiras de carência e luxúria até à
paixão justa de «querer matar e aniquilar estes degenerados!»

Contempla isso como algo dentro da tua mente. Já vi ambas as tendências em mim
mesmo. Posso sentir-me atraído por prazeres sensuais. E também posso ser muito
duro e hipócrita, e crítico dos outros e de mim mesmo. O \emph{bhava-taṇhā}
também pode ser muito doce quando o fazemos pelo bem-estar dos outros. Não é
apenas que eu queira alcançar algo para poder dizer que alcancei algo. Há também
o \emph{bhava-taṇhā} de querer porque sentimos que gostaríamos de ajudar todos
os outros. Continua a existir o «eu sou». «Quero alcançar a iluminação e,
depois, vou ajudar verdadeiramente todos os outros e quero tornar-me alguém que
não é egoísta, mas que trabalha totalmente pelo bem-estar de todos os seres
sencientes.» Isso é muito altruísta, não é? É belo e inspirador, mas também pode
ser \emph{bhava-taṇhā} se provir de \emph{avijjā-paccayā saṅkhārā.}

Quando vemos com clareza através do \emph{vijjā} e contemplamos o Dhamma, então
não há ninguém que se torne algo, que conquiste ou que alcance. As coisas são
como são. O bem é feito e o mal é evitado nas ações e nas palavras. Há o ato de
fazer o bem. O que resta para fazer na vida senão ser virtuoso? Não é essa a
beleza da nossa humanidade? O que é verdadeiramente alegre e encantador em ser
humano é a nossa capacidade de ser virtuoso. A experiência humana destina-se à
virtude e à bondade e a abster-se de fazer coisas más e prejudiciais a nós
próprios e aos outros. Não consigo pensar em mais nada que valha a pena fazer!

\photoPage{image-013-bhikkhu-chainsaw.png}{Bhikkhu a serrar troncos}{}
